\section{Introduction}
\label{sec:introduction}

\IEEEPARstart{T}{he} development of foundation models and distributed intelligent systems has accelerated the transition in artificial intelligence from passive sequence-to-sequence text generation toward autonomous, trustworthy \textit{Agentic AI} architectures~\cite{min2023recent, wei2024trustworthy, chander2024toward, goyal2023survey, standen2024adversarial, dehghantanha2026sok, deng2025ai}. Contemporary agentic architectures synthesize Large Language Models (LLMs) with algorithmic planning engines, long-term state repositories, external knowledge retrievers, and execution interfaces (such as tools, software APIs, and operating system shells)~\cite{narajala2025securing}. By executing iterative loops of perception, cognitive reasoning, tool invocation, and environment feedback, these agents perform multi-horizon workflows across software engineering, autonomous data analysis, enterprise resource planning, and cyber operations~\cite{liu2024demystifying, helpnet2023genai}.

Despite these operational capabilities, the transition toward autonomous agency introduces fundamental dependability and security vulnerabilities~\cite{iyer2007editorial, schneier2023trustworthy}. In traditional computational systems, execution environments maintain strict hardware- and kernel-enforced isolation between untrusted data streams and executable control logic~\cite{dehghantanha2026sok}. In contrast, LLM-based agentic architectures process natural language instructions, operational guidelines, retrieved third-party content, conversation history, and tool definitions within a single, shared semantic context space, creating vulnerability to adversarial manipulations~\cite{greshake2023not, sutton2024indirect}. Consequently, probabilistic foundation models cannot deterministically segregate authoritative system instructions from adversarial payloads embedded in external data~\cite{chen2025struq, yi2025benchmarking, vassilev2026robust}.

When an LLM is granted effector privileges (such as executing shell scripts, issuing HTTP requests, mutating database records, or initiating financial transactions), the failure to isolate code from data creates severe systemic risks~\cite{dehghantanha2026sok, liu2024demystifying}. Empirical investigations demonstrate that prompt manipulations can trigger Remote Code Execution (RCE)~\cite{liu2024demystifying}, targeted data poisoning can subvert Retrieval-Augmented Generation (RAG) pipelines~\cite{zou2025poisonedrag, an2025rag, shafran2025machine}, and malicious cross-agent signaling can provoke cascading failures across multi-agent systems~\cite{standen2024adversarial, gu2024agent, peigne2025multi, chen2024blockagents}.

\subsection{Foundations and Limitations of Prior Systematizations}
\label{subsec:prior_work_limits}

In a foundational Systematization of Knowledge (SoK), Dehghantanha and Homayoun~\cite{dehghantanha2026sok} mapped the primary trust boundaries and attack surfaces of tool-augmented LLM systems. Their work codified five canonical threat paths (P1--P5), defined seven attacker-aware evaluation metrics (such as Unsafe Action Rate and Privilege Escalation Distance), and synthesized early evidence from the 2023--2025 research epoch. Other contemporary overviews examine specific facets of generative security and trustworthy distributed intelligence: Min~\textit{et al.}~\cite{min2023recent} reviewed foundational LLM developments; Wei and Liu~\cite{wei2024trustworthy} analyzed robustness, privacy, and governance in distributed AI; Chander~\textit{et al.}~\cite{chander2024toward} surveyed explainability and robustness in trustworthy AI; Goyal~\textit{et al.}~\cite{goyal2023survey} categorized adversarial defenses in NLP; and Standen~\textit{et al.}~\cite{standen2024adversarial} investigated adversarial machine learning in multi-agent environments; alongside agent vulnerability surveys by Deng~\textit{et al.}~\cite{deng2025ai}; high-level threat models by Narajala and Narayan~\cite{narajala2025securing}; trust, risk, and security governance (TRiSM) by Raza~\textit{et al.}~\cite{raza2026trism}; and coordination challenges by Schroeder de Witt~\textit{et al.}~\cite{dewitt2025open}. 

However, existing literature exhibits five structural limitations:
\begin{enumerate}
    \item \textbf{Omission of Multimodal and Computer-Use Frontiers:} Prior surveys focus almost exclusively on text-based function calling, leaving graphical user interface (GUI) automation, visual prompt injection, and OS-level Computer-Use architectures underexplored~\cite{gu2024agent, xie2024osworld}.
    \item \textbf{Absence of Standardized Tool Protocol Analysis:} The rapid standardization of agent interfaces, exemplified by the Model Context Protocol (MCP)~\cite{hou2026mcp} and OpenAI function calling standards, introduces novel host-client-server trust dynamics and tool shadowing risks that have not yet been systematically formalized.
    \item \textbf{Superficial Multi-Agent System (MAS) Topologies:} While initial studies acknowledge inter-agent prompt propagation~\cite{gu2024agent, dewitt2025open}, they do not systematically categorize vulnerabilities arising from heterogeneous coordination topologies (such as hierarchical orchestrators, peer-to-peer swarms, blackboard architectures, and consensus voting mechanisms).
    \item \textbf{Bifurcation of Empirical Guardrails and Formal Verification:} The literature remains divided between heuristic, empirical prompt filters~\cite{inan2023llama, zeng2024shieldgemma} and formal systems security principles~\cite{chen2025struq}. A unified framework integrating capability-based access control (CBAC), temporal logic plan verification, and runtime invariant enforcement is conspicuously lacking.
    \item \textbf{Fragmented Benchmarking Testbeds:} Current security benchmarks (such as AgentDojo~\cite{debenedetti2024agentdojo}, InjecAgent~\cite{zhan2024injecagent}, and ToolEmu~\cite{ruan2023identifying}) operate in disjoint evaluation silos with divergent threat assumptions, hindering reproducible cross-paradigm comparisons.
\end{enumerate}

\subsection{Contributions and Organization}
\label{subsec:contributions}

To bridge these gaps, this paper establishes a systematic survey that broadens and formalizes the security analysis of agentic AI. The primary contributions of this work are as follows:
\begin{itemize}
    \item \textbf{Comprehensive Architectural Decomposition:} We formulate an end-to-end reference architecture encompassing the full perception-cognition-action-reflection lifecycle, systematically mapping the expanded Trusted Computing Base (TCB) across single-agent, hierarchical, and decentralized coordination topologies.
    \item \textbf{Granular Attack Surface Taxonomy:} We establish an analytical taxonomy of exploitation vectors, decomposing attack goals (G1--G7) into concrete micro-primitives, including memory sleep-cycle corruption, semantic cache poisoning, reflection hijacking, tool parameter smuggling, and Insecure Direct Object Reference (IDOR) vulnerabilities in retrieval embeddings.
    \item \textbf{Frontier Analysis of Protocols and Modalities:} We provide a systematic security treatment of the Model Context Protocol (MCP) and investigate multimodal attack surfaces, including visual typographic injection and OS-level GUI manipulation.
    \item \textbf{Multi-Agent Threat Modeling:} We categorize topology-specific vulnerabilities in Multi-Agent Systems (MAS), modeling Byzantine consensus degradation, collective hallucination loops, and self-replicating prompt worms.
    \item \textbf{Holistic Defense-in-Depth Framework:} We formulate a tiered defensive taxonomy that bridges empirical guardrails, structured query enforcement, capability-based sandboxing, and formal contract-based plan verification.
    \item \textbf{Unified Benchmark and Governance Synthesis:} We present a comparative matrix of contemporary agent security benchmarks and map technical vulnerabilities to emerging international governance frameworks (OWASP GenAI Top 10, MITRE ATLAS, NIST AI RMF, and the EU AI Act).
\end{itemize}

The remainder of this paper is structured as follows. Section~\ref{sec:architecture} introduces the agentic AI reference architecture and operational lifecycle. Section~\ref{sec:threat_model} formalizes the threat model, adversary profiles, and causal threat graph. Section~\ref{sec:attack_taxonomy} establishes the unified taxonomy of attack vectors and multi-stage kill chains. Section~\ref{sec:modality_mcp} examines multimodal GUI and protocol-level vulnerabilities. Section~\ref{sec:multi_agent} investigates multi-agent coordination risks. Section~\ref{sec:defenses} synthesizes multi-layer defensive paradigms. Section~\ref{sec:benchmarks} evaluates security metrics and benchmarking frameworks. Section~\ref{sec:standards} analyzes governance and standards. Section~\ref{sec:open_challenges} articulates open research problems, and Section~\ref{sec:conclusion} concludes the paper.
